\documentclass[11pt]{article}

\usepackage[margin=1in]{geometry}
\usepackage{amsmath,amssymb,amsthm,mathtools}
\usepackage[hidelinks]{hyperref}

\newtheorem{theorem}{Theorem}[section]
\newtheorem{lemma}[theorem]{Lemma}
\newtheorem{proposition}[theorem]{Proposition}
\theoremstyle{remark}
\newtheorem{remark}[theorem]{Remark}

\newcommand{\Arg}{\operatorname{Arg}}
\newcommand{\tr}{\operatorname{tr}}

\title{Positive Square Energy for Two Pentagons Joined by an Arbitrary Path}
\author{}
\date{25 July 2026}

\begin{document}
\maketitle

\begin{abstract}
We prove a uniform positive-square-energy bound for every connected cactus
whose only cyclic blocks are two vertex-disjoint pentagons.  The pentagons may
be joined by a path of any positive length, and arbitrary trees may be attached
anywhere.  If the graph has order $n$, then
\[
 s^+(G)>n+5-2\sqrt5>n.
\]
The proof does not compare the graph with its bare bicyclic core: attached
trees can increase that core's pointwise Coulson phase.  Instead, the joining
path is factored exactly, and the phase of the joined graph is compared with
the sum of the phases of the two actual unicyclic lobes.
\end{abstract}

\section{Statement and notation}

All graphs are finite, simple, and undirected.  If the adjacency eigenvalues
of a graph $H$ are $\lambda_1,\ldots,\lambda_h$, put
\[
 s^+(H)=\sum_{\lambda_j>0}\lambda_j^2,
 \qquad
 s^-(H)=\sum_{\lambda_j<0}\lambda_j^2,
 \qquad
 D(H)=s^+(H)-s^-(H).
\]

\begin{theorem}\label{thm:main}
Let $G$ be a connected cactus whose only cyclic blocks are two copies of
$C_5$.  Suppose that the two cycles are vertex-disjoint and that the unique
path joining them has positive length.  Arbitrary trees may be attached at
any vertices of the cycles or the joining path.  If $n=|V(G)|$, then
\begin{equation}\label{eq:main}
 s^+(G)>n+5-2\sqrt5>n.
\end{equation}
\end{theorem}

Thus the theorem includes a single bridge between the cycles and every longer
joining path.  It does not include the bouquet in which the two pentagons
share a vertex; that graph has no positive-length path with distinct cycle
endpoints and requires a separate theorem.

\section{Normalized characteristic polynomials}

For a graph $F$, define its signless matching partition by
\begin{equation}\label{eq:Z-unweighted}
 Z_F(t)=\sum_M t^{|V(F)|-2|M|},
\end{equation}
where the sum is over all matchings.  More generally, for positive vertex
activities $a=(a_v)$, write
\begin{equation}\label{eq:Z-weighted}
 Z_F(a)=\sum_M\prod_{v\text{ unmatched by }M}a_v.
\end{equation}
All edge activities are one.  Every such partition is positive.

For a graph $H$ of order $h$, let
\begin{equation}\label{eq:Psi}
 \Psi_H(t)=i^{-h}\det(itI-A(H))
          =\prod_{j=1}^h(t+i\lambda_j),
 \qquad t>0.
\end{equation}
Grouping Sachs' expansion by its collection $\mathcal C$ of pairwise
vertex-disjoint cycle components gives
\begin{equation}\label{eq:Sachs}
 \Psi_H(t)=\sum_{\mathcal C}
 \prod_{C\in\mathcal C}(-2i^{-|C|})
 Z_{H-V(\mathcal C)}(t).
\end{equation}
Indeed, after the cycles are fixed, the remaining elementary components are
a matching; substitution of $it$ changes its alternating matching polynomial
into~\eqref{eq:Z-unweighted}.  In particular, a pentagon has multiplier
\begin{equation}\label{eq:C5-multiplier}
 -2i^{-5}=2i.
\end{equation}

We use the continuous argument
\begin{equation}\label{eq:continuous-argument}
 \Theta_H(t)=\sum_{j=1}^h\arctan\frac{\lambda_j}{t},
 \qquad -\frac\pi2<\arctan x<\frac\pi2.
\end{equation}
It is the continuous argument of~\eqref{eq:Psi} that tends to zero as
$t\to\infty$.

\begin{lemma}[Signed Coulson identity]\label{lem:coulson}
For every graph $H$,
\begin{equation}\label{eq:Coulson}
 D(H)=-\frac4\pi\int_0^\infty t\Theta_H(t)\,dt.
\end{equation}
The integral is absolutely convergent.
\end{lemma}

\begin{proof}
For every real $\lambda$,
\[
 \lambda|\lambda|=\frac2\pi\int_0^\infty
 \frac{\lambda^3}{\lambda^2+t^2}\,dt.
\]
Also
\[
 \Theta_H'(t)=-\sum_j\frac{\lambda_j}{\lambda_j^2+t^2},
 \qquad
 \sum_j\frac{\lambda_j^3}{\lambda_j^2+t^2}
 =-t^2\sum_j\frac{\lambda_j}{\lambda_j^2+t^2},
\]
where $\sum_j\lambda_j=\tr A(H)=0$ was used in the second identity.
Summing the scalar formula and integrating by parts therefore gives
\eqref{eq:Coulson}.  At zero, $t^2\Theta_H(t)\to0$.  At infinity,
$\sum_j\lambda_j=0$ and the arctangent expansion give
$\Theta_H(t)=O(t^{-3})$.  These estimates also prove absolute convergence.
\end{proof}

\section{Tree splitting and the two actual lobes}

Let the joining path have endpoints $u_1$ and $u_2$ on the two pentagons and
let it have $m\geq0$ internal vertices.  Thus $m=0$ is the bridge case.  Put
into $X_i$ the $i$th pentagon together with every tree branch attached at its
cycle vertices, and root $X_i$ at $u_i$.  The graphs $X_1,X_2$ are the
\emph{actual} pentagonal unicyclic lobes; no cycle attachment is discarded.
Only tree branches attached to internal vertices of the joining path remain
to be eliminated.

We record the elementary matching belief-propagation split.  Orient a rooted
tree branch toward the path.  If $x\to p$ is an oriented edge and $T_{x\to p}$
is the component below $x$, including $x$, set
\begin{equation}\label{eq:message}
 q_{x\to p}(t)=\frac{Z_{T_{x\to p}}(t)}{Z_{T_{x\to p}-x}(t)}.
\end{equation}
Splitting a matching according to the status of $x$ gives
\begin{equation}\label{eq:BP}
 q_{x\to p}(t)=t+
 \sum_{y\text{ child of }x}\frac1{q_{y\to x}(t)}\geq t>0.
\end{equation}
After all path-internal branches are split off, the $j$th internal path
vertex has a positive effective activity
\begin{equation}\label{eq:activity}
 a_j(t)=t+\sum_{x\text{ attached at }j}\frac1{q_{x\to j}(t)}>0,
\end{equation}
and all terms have a common positive real prefactor
\begin{equation}\label{eq:B0}
 B_0(t)=\prod_{x\to j\text{ at an internal path vertex}}
 Z_{T_{x\to j}}(t)>0.
\end{equation}
To verify the common factor even when a vertex is removed by a matching edge,
note that extracting $Z_T$ gives the edge into a branch the weight
$Z_{T-x}/Z_T=1/q_{x\to j}$.  If the path vertex is not retained, the detached
branch contributes its full factor $Z_T$.  Thus no hidden quotient remains.

Define
\begin{equation}\label{eq:zwr}
 z_i=\Psi_{X_i}(t),
 \qquad
 w_i=\Psi_{X_i-u_i}(t),
 \qquad
 r_i=\frac{w_i}{z_i}.
\end{equation}
The graph $X_i-u_i$ is a forest, so $w_i=Z_{X_i-u_i}(t)>0$.  Formula
\eqref{eq:Sachs} applied to the unique pentagon of $X_i$ gives
\begin{equation}\label{eq:lobe-quadrant}
 z_i=Z_{X_i}(t)+2iZ_{X_i-V(C_5)}(t).
\end{equation}
Both displayed matching partitions are positive.  Hence $z_i$ is in the
open first quadrant.  Writing
\begin{equation}\label{eq:r-coordinates}
 r_i=x_i-iy_i
\end{equation}
therefore gives $x_i>0$ and $y_i>0$: each $r_i$ is in the open fourth
quadrant.  Put $\theta_i=\Arg z_i\in(0,\pi/2)$.

\section{The exact connector factor}

For a list of variables define the continuant by
\begin{equation}\label{eq:continuant}
 K()=1,\qquad K(b_1)=b_1,\qquad
 K(b_1,\ldots,b_k)=b_kK(b_1,\ldots,b_{k-1})
                    +K(b_1,\ldots,b_{k-2}).
\end{equation}

\begin{proposition}[Connector factor]\label{prop:connector}
With the preceding notation,
\begin{equation}\label{eq:factorization}
 \Psi_G(t)=B_0(t)z_1z_2F(t),
\end{equation}
where
\begin{equation}\label{eq:F-cases}
 F=\begin{cases}
 1+r_1r_2,&m=0,\\
 a_1+r_1+r_2,&m=1,\\
 K(a_1+r_1,a_2,\ldots,a_{m-1},a_m+r_2),&m\geq2.
 \end{cases}
\end{equation}
Moreover, $\operatorname{Im}F<0$.
\end{proposition}

\begin{proof}
First ignore the already extracted factor $B_0$.  In each grouped Sachs term,
split matchings along the joining path.  Summing over the two possible Sachs
states inside a lobe gives $z_i$ when its boundary edge is unused; using that
edge deletes the root and gives $w_i$.  Thus $w_i/z_i=r_i$ records use of the
boundary edge at the $i$th lobe.  For $m=0$, the bridge is either unused or used, giving
$1+r_1r_2$.  For $m=1$, the internal vertex is unmatched, or is matched to
exactly one boundary, giving $a_1+r_1+r_2$.  For $m\geq2$, the usual last-edge
split for path matchings is exactly recurrence~\eqref{eq:continuant}, with
the two boundary contributions added to the first and last activities.  This
proves~\eqref{eq:factorization}--\eqref{eq:F-cases}.

For the sign of the imaginary part, expand
\begin{equation}\label{eq:ABCD}
 F=A+Br_1+Cr_2+Dr_1r_2.
\end{equation}
Here $B$ is the coefficient of $r_1$ and $C$ is the coefficient of $r_2$.
For $m=0$ one has $(A,B,C,D)=(1,0,0,1)$, while for $m=1$ one has
$(A,B,C,D)=(a_1,1,1,0)$.  If $m\geq2$, the matching expansion of the
continuant shows that all four coefficients are positive: $A,B,C,D$ are the
positive partition sums obtained by prescribing neither, only the left,
only the right, or both boundary edges, respectively.  These four classes
are nonempty for every $m\geq2$.

Using~\eqref{eq:r-coordinates},
\begin{equation}\label{eq:ImF}
 \operatorname{Im}F
 =-By_1-Cy_2-D(x_1y_2+x_2y_1)<0.
\end{equation}
For $m=0$ the last term is strictly negative; for $m=1$ the middle two terms
are strictly negative; and for $m\geq2$ all three displayed contributions
are negative.  This proves the assertion in every case.
\end{proof}

The identities in Proposition~\ref{prop:connector} are independently
corroborated by the exact symbolic and random-cactus audit script
\begin{center}
\texttt{positive-square-energy/experiments/connector\_factor\_audit.py}.
\end{center}
The script is not needed for the proof.

\section{Strict phase comparison}

\begin{proposition}\label{prop:phase}
For every $t>0$,
\begin{equation}\label{eq:phase-bound}
 0<\Theta_G(t)<\theta_1(t)+\theta_2(t).
\end{equation}
Consequently,
\begin{equation}\label{eq:D-split}
 D(G)>D(X_1)+D(X_2).
\end{equation}
\end{proposition}

\begin{proof}
Because the only cycles of $G$ are its two pentagons, the imaginary part of
the grouped Sachs expansion~\eqref{eq:Sachs} is exactly
\begin{equation}\label{eq:whole-upper}
 \operatorname{Im}\Psi_G(t)
 =2Z_{G-V(C_5^{(1)})}(t)+2Z_{G-V(C_5^{(2)})}(t)>0.
\end{equation}
The empty-cycle and double-cycle terms are real.  Thus $\Psi_G(t)$ lies in
the open upper half-plane, and its continuous argument satisfies
$0<\Theta_G(t)<\pi$.

By Proposition~\ref{prop:connector}, the principal argument
$\varphi=\Arg F$ belongs to $(-\pi,0)$.  Since
$\theta_1+\theta_2\in(0,\pi)$, the number
$\theta_1+\theta_2+\varphi$ lies in $(-\pi,\pi)$ and is congruent modulo
$2\pi$ to $\Theta_G$.  It cannot be nonpositive: adding $2\pi$ would then
put the corresponding argument in $(\pi,2\pi)$, contradicting
$0<\Theta_G<\pi$.  Therefore
\begin{equation}\label{eq:branch-identity}
 \Theta_G=\theta_1+\theta_2+\Arg F,
\end{equation}
and~\eqref{eq:phase-bound} follows strictly from $\Arg F<0$.

Apply the Coulson identity to $G,X_1,X_2$.  All three phase integrals converge,
and the integrand
$t(\theta_1+\theta_2-\Theta_G)$ is continuous and strictly positive on
$(0,\infty)$.  Hence
\[
 D(G)-D(X_1)-D(X_2)
 =\frac4\pi\int_0^\infty
 t(\theta_1+\theta_2-\Theta_G)\,dt>0,
\]
which is~\eqref{eq:D-split}.
\end{proof}

This is deliberately not a comparison with the phase of the bare two-cycle
core.  Such a comparison is false in general: attachments may increase the
bare-core phase pointwise.  The valid comparison is with the sum of the
phases of the actual lobes $X_1$ and $X_2$, after which the connector supplies
the strictly negative phase $\Arg F$.

\section{The pentagonal lobe bound}

\begin{lemma}\label{lem:lobe}
If $X$ is a unicyclic graph whose unique cycle is $C_5$, then
\begin{equation}\label{eq:lobe-D}
 D(X)\geq D(C_5)=4-2\sqrt5.
\end{equation}
\end{lemma}

\begin{proof}
Apply the same branch split~\eqref{eq:message}--\eqref{eq:BP} to all trees
attached to the five cycle vertices.  It gives a common positive factor and
effective cycle activities $b_v=t+y_v\geq t$.  Formula~\eqref{eq:Sachs}
becomes
\begin{equation}\label{eq:weighted-cycle}
 \Psi_X(t)=K_0(t)\bigl(Z_{C_5}(b)+2i\bigr),
 \qquad K_0(t)>0.
\end{equation}
The weighted matching partition has nonnegative coefficients in all vertex
activities, so
\begin{equation}\label{eq:subpartition}
 Z_{C_5}(b)\geq Z_{C_5}(t).
\end{equation}
Both quantities in~\eqref{eq:weighted-cycle} lie in the first quadrant.
Consequently
\[
 0<\Theta_X(t)=\arctan\frac2{Z_{C_5}(b)}
 \leq\arctan\frac2{Z_{C_5}(t)}=\Theta_{C_5}(t).
\]
The Coulson identity yields $D(X)\geq D(C_5)$.

For completeness, the eigenvalues of $C_5$ are
$2\cos(2\pi j/5)$ for $0\leq j<5$.  Its positive eigenvalues are $2$ and
$(\sqrt5-1)/2$ with multiplicity two; hence
\[
 s^+(C_5)=4+2\left(\frac{\sqrt5-1}{2}\right)^2
          =7-\sqrt5.
\]
Since $C_5$ has five edges,
$s^+(C_5)+s^-(C_5)=\tr A(C_5)^2=10$, and therefore
$D(C_5)=2s^+(C_5)-10=4-2\sqrt5$.
\end{proof}

\begin{proof}[Proof of Theorem~\ref{thm:main}]
Proposition~\ref{prop:phase} and Lemma~\ref{lem:lobe} give
\begin{equation}\label{eq:final-D}
 D(G)>2(4-2\sqrt5)=8-4\sqrt5.
\end{equation}
The connected graph $G$ has cycle rank two, so if $e=|E(G)|$, then
$e=n+1$.  Therefore
\[
 s^+(G)+s^-(G)=\tr A(G)^2=2e=2n+2.
\]
Averaging this identity with~\eqref{eq:final-D} proves
\[
 s^+(G)=n+1+\frac{D(G)}2
 >n+1+4-2\sqrt5=n+5-2\sqrt5.
\]
Finally, $5-2\sqrt5>0$ because $25>20$, completing the proof.
\end{proof}

\end{document}
